\section{Related work}
\label{sec:related}
\paragraph{Cheaper evaluation inside a benchmark.}
tinyBenchmarks \citep{polo2024tiny} and metabench \citep{kipnis2024metabench} shrink the number of \emph{items} used to estimate a model's score on a given benchmark; metabench selects items with item-response-theory (IRT) information, reconstructs the scores of six benchmarks from a sparse item set, and compares against random subsets of equal size. We borrow the random-subset null and the ``reconstruct what was not run'' framing, but our unit of selection is the whole benchmark: the evaluator picks which benchmarks to run, and the remaining benchmarks are \emph{not} run at all. Item-level reduction and benchmark-level reduction are complementary.

\paragraph{Low-dimensional structure of model capabilities.}
Observational scaling laws \citep{ruan2024observational} extract a few principal ``capability'' dimensions from a benchmark$\times$model table and use them to predict downstream metrics; Sloth \citep{polo2024sloth} fits low-dimensional skill models with family-specific scaling and needs model-side covariates (size, tokens) and at least one model of the target family. We use the same structural idea (a low-rank latent space) but in a cold-start setting: the covariates of the new model are its \emph{measured} scores on $k$ benchmarks, no model metadata is used, and the question is which $k$ benchmarks to measure and how many. \citet{perlitz2024benchbench} show that benchmark agreement conclusions depend on methodological choices; our protocol therefore states its choices (splits, targets, normalisation) and varies them.

\paragraph{Subset selection and optimal design.}
Choosing columns of a matrix that best represent the rest is the column subset selection / CUR problem \citep{mahoney2009cur}, with leverage-score and pivoted-QR heuristics. Under a Gaussian model, choosing observations that minimise the posterior variance of unobserved quantities is a classical sensor-placement problem; greedy mutual-information or log-determinant selection has a $(1-1/e)$ guarantee for monotone submodular objectives \citep{nemhauser1978analysis,guestrin2005near}, and spectral conditions govern greedy variance reduction \citep{das2011submodular}. Randomly pivoted Cholesky \citep{chen2025rpcholesky} is a recent relative. We use these as selectors and check, rather than assume, that they help on real score matrices.

\paragraph{Missing data and cold start.}
Our inference rules build on factor analysis and PCA with missing values \citep{ilin2010pca}, matrix completion \citep{candes2009exact} and matrix-factorisation recommenders \citep{koren2009mf}. The new model corresponds to a cold-start user, the benchmarks to items, and the $k$ benchmarks to the ``interview items'' of the item selection problem for cold-start recommendation \citep{meng2020item}. Missingness in leaderboards is not at random in Rubin's sense \citep{rubin1976missing} (benchmarks enter a leaderboard with its version, models with their release date), which we address with held-out-family and time-split evaluations. Psychometric treatments of leaderboards \citep{federiakin2025psychometric} are also related.

\paragraph{Benchmark suites used here.}
MTEB \citep{muennighoff2022mteb} is our main data; HELM \citep{liang2022helm}, the Open LLM Leaderboard and LiveBench provide small generative-LLM cross-checks (Section~\ref{sec:data}). MMLU \citep{hendrycks2020mmlu} underlies the HELM-MMLU subject matrix.
